% 第 7 章：频率-持续时间解析法
\section{频率-持续时间解析法}\label{sec:rel-fd}

\subsection{容量停运概率表（COPT）}
\srcpath{run\_frequency\_duration\_analysis}（\srcpath{reliability\_assessment.hpp:666}）不做
仿真，递归构建 COPT 后解析给出 LOLP/LOLE/LOLF，假设元件故障相互独立。逐台加入一个
容量 $C$、不可用度 $U$ 的两态机组时，停运容量分布按卷积更新
\begin{equation}
P^{new}(X)=(1-U)\,P^{old}(X)+U\,P^{old}(X-C),
\end{equation}
并同时维护累积概率 $P(\text{停运}\ge X)$ 与累积频率 $F(\text{停运}\ge X)$（占/年），
结果字段 \fld{capacity\_outage\_levels}/\fld{cumulative\_probability}/\fld{cumulative\_frequency}
（\fld{FrequencyDurationResult}，:224）。

\subsection{系统指标}
设峰荷 $L$（\srcpath{peak\_load\_mw}，缺省取母线负荷之和）、总装机 $C_{tot}$，则失负荷发生于
可用容量小于负荷，即停运超过备用 $R=C_{tot}-L$：
\begin{equation}
\mathrm{LOLP}=P(\text{停运}>R),\quad
\mathrm{LOLE}_{fd}=\mathrm{LOLP}\cdot H,\quad
\mathrm{LOLF}_{fd}=F(\text{停运}>R),
\end{equation}
\begin{equation}
\mathrm{LOLD}=\frac{\mathrm{LOLE}_{fd}}{\mathrm{LOLF}_{fd}}\quad(\text{小时/次}),
\end{equation}
分别对应 \fld{lolp}、\fld{lole\_fd}、\fld{lolf\_fd}、\fld{lold}。

\subsection{适用边界}
本法比 MC 快且无采样方差，但\textbf{假设元件故障独立}、以两态机组容量模型为主，不含
时序负荷、修复时序、共因与网络约束（切负荷由容量裕度而非潮流决定）。需要网络受限的
失负荷或时序尾部风险时应改用 FMEA 或序贯 MC。

\subsection{配电可靠性指标与尾部风险}
\srcpath{compute\_distribution\_indices}（:679）依 IEEE Std 1366-2012 由逐母线中断频率
\fld{nodal\_cif} 与时长 \fld{nodal\_cid}、客户数 \srcpath{Load::num\_customers} 聚合
\fld{saifi}/\fld{saidi}/\fld{caidi}/\fld{asai}/\fld{asui}（\fld{DistributionIndices}，:208）。
DC 感知：当逐点向量长度超过 AC 母线数时追加 DC 母线客户，布局为 \texttt{[AC | DC]}。
尾部风险 \srcpath{compute\_tail\_risk}（:672）给 EENS/LOLE 的 VaR/CVaR 与分位数
（\fld{TailRiskMetrics}，:187）。
